\documentclass[11pt,a4paper]{article}
\usepackage[utf8]{inputenc}
\usepackage[vietnamese]{babel}
\usepackage{amsmath,amssymb}
\usepackage{geometry}
\geometry{a4paper, total={180mm,260mm}, left=15mm, right=15mm, top=18mm, bottom=18mm}
\usepackage{tikz}
\usetikzlibrary{shapes.geometric, arrows.meta, positioning, calc}
\usepackage{tcolorbox}
\tcbuselibrary{skins,breakable}
\usepackage{xcolor}
\usepackage{multicol}
\usepackage{array}
\usepackage{fontawesome}
\usepackage{fancyhdr}

\pagestyle{fancy}
\fancyhf{}
\renewcommand{\headrulewidth}{0.5pt}
\renewcommand{\footrulewidth}{0.5pt}
\fancyhead[L]{\small\bfseries\color{teal}\faBook\ SỔ TAY KIẾN THỨC TOÁN 3}
\fancyhead[R]{\small\bfseries\color{teal}Lớp Toán Thầy Lĩnh -- loptoanthaylinh.vn}
\fancyfoot[C]{\small\color{gray}-- Sổ tay kiến thức: Bài 37. Chia số có ba chữ số cho số có một chữ số --}

\definecolor{primary}{HTML}{005F73}
\definecolor{secondary}{HTML}{0A9396}
\definecolor{accent}{HTML}{EE9B00}
\definecolor{danger}{HTML}{AE2012}
\definecolor{boxbg}{HTML}{F4F6F8}

\begin{document}

% --- BANNER TIÊU ĐỀ ---
\begin{tcolorbox}[
  enhanced,
  colback=white,
  colframe=primary,
  arc=3mm,
  boxrule=1.2pt,
  left=4mm, right=4mm, top=2.5mm, bottom=2.5mm,
  drop shadow=black!10!white
]
  \centering
  {\Large\bfseries\color{primary}\faBookmark\ BÀI 37: CHIA SỐ CÓ BA CHỮ SỐ CHO SỐ CÓ MỘT CHỮ SỐ}\\[1.5mm]
  {\small\itshape\color{black!75}Quy trình Chia - Nhân - Trừ từ trái sang phải, chia hết và phép chia có chữ số 0 ở thương trên 1 trang A4}
\end{tcolorbox}

\vspace{-2mm}


\noindent
\begin{minipage}[t]{0.49\textwidth}\vspace{0pt}
\begin{tcolorbox}[
  enhanced,
  colback=blue!4!white,
  colframe=blue!75!black,
  arc=2mm, boxrule=0.8pt,
  left=2mm, right=2mm, top=1.8mm, bottom=1.8mm,
  title={\footnotesize\bfseries\color{white}\faCalculator\ 1. Quy trình đặt tính và chia}
]
{\footnotesize
\begin{itemize}\setlength{\itemsep}{1.5pt}\setlength{\leftskip}{-2.5mm}
  \item Thực hiện chia theo thứ tự từ \textbf{trái sang phải} (từ hàng trăm đến hàng đơn vị).
  \item Mỗi lượt chia thực hiện theo nhịp 3 bước:
  $$\mathbf{\text{CHIA}} \implies \mathbf{\text{NHÂN}} \implies \mathbf{\text{TRỪ}}$$
  \item \textbf{Lượt chia đầu tiên:} Nếu chữ số hàng trăm nhỏ hơn số chia thì phải lấy cả 2 chữ số đầu để chia.
  \item \textbf{Từ lượt chia thứ hai:} Mỗi lần chỉ hạ đúng $\mathbf{1}$ chữ số.
\end{itemize}
}
\end{tcolorbox}
\end{minipage}
\hfill
\begin{minipage}[t]{0.49\textwidth}\vspace{0pt}
\begin{tcolorbox}[
  enhanced,
  colback=teal!4!white,
  colframe=teal!75!black,
  arc=2mm, boxrule=0.8pt,
  left=2mm, right=2mm, top=1.8mm, bottom=1.8mm,
  title={\footnotesize\bfseries\color{white}\faKey\ 2. Trường hợp đặc biệt chữ số 0 ở thương}
]
{\footnotesize
\begin{itemize}\setlength{\itemsep}{1.5pt}\setlength{\leftskip}{-2.5mm}
  \item Ở lượt chia bất kì sau lượt đầu, nếu số bị chia nhỏ hơn số chia:
  $$\implies \text{Viết } \mathbf{0} \text{ vào thương}$$
  rồi mới hạ chữ số tiếp theo!
  \item \textbf{Kiểm tra tính đúng đắn:}
  $$\text{Số bị chia} = \text{Thương} \times \text{Số chia} + \text{Số dư}$$
  Số dư luôn phải \textbf{nhỏ hơn số chia} ($r < d$).
\end{itemize}
}
\end{tcolorbox}
\end{minipage}

\vspace{-2.5mm}
% --- KHỐI 3: SƠ ĐỒ TƯ DUY TIKZ ---
\noindent
\begin{tcolorbox}[
  enhanced,
  colback=orange!4!white,
  colframe=orange!80!black,
  arc=2mm, boxrule=0.8pt,
  left=2mm, right=2mm, top=1.5mm, bottom=1.5mm,
  title={\footnotesize\bfseries\color{white}\faSitemap\ SƠ ĐỒ TƯ DUY: VÒNG LẶP 3 BƯỚC ĐẶT TÍNH CHIA}
]
\centering
\begin{tikzpicture}[
  box/.style={draw, rounded corners=2pt, align=center, font=\scriptsize\bfseries, inner sep=2.5pt, minimum height=0.75cm},
  arr/.style={-{Stealth[length=2mm]}, thick}
]
  \node[box, fill=purple!15, draw=purple!80!black, text width=2.6cm] (c1) at (0, 0) {BƯỚC 1: CHIA\\Ước lượng thương\\Ghi chữ số vào thương};
  \node[box, fill=teal!15, draw=teal!80!black, text width=2.6cm] (c2) at (3.5, 0) {BƯỚC 2: NHÂN\\Nhân thương vừa tìm\\với số chia};
  \node[box, fill=blue!12, draw=blue!80!black, text width=2.6cm] (c3) at (7.0, 0) {BƯỚC 3: TRỪ\\Trừ để tìm số dư\\Số dư phải $< d$};
  \node[box, fill=orange!15, draw=orange!80!black, text width=2.6cm] (c4) at (10.5, 0) {HẠ TIẾP\\Hạ 1 chữ số tiếp theo\\Lặp lại vòng 3 bước};

  \draw[arr, purple!80!black] (c1) -- (c2);
  \draw[arr, teal!80!black] (c2) -- (c3);
  \draw[arr, blue!80!black] (c3) -- (c4);
\end{tikzpicture}
\end{tcolorbox}

\vspace{-2.5mm}
% --- KHỐI 4: VÍ DỤ & MẸO TRÁNH BẪY ---
\noindent
\begin{minipage}[t]{0.49\textwidth}\vspace{0pt}
\begin{tcolorbox}[
  enhanced,
  colback=green!4!white,
  colframe=green!60!black,
  arc=2mm, boxrule=0.8pt,
  left=2mm, right=2mm, top=1.8mm, bottom=1.8mm,
  title={\footnotesize\bfseries\color{white}\faCheckCircle\ Ví dụ thực hành chuẩn mực}
]
{\footnotesize
\textbf{1. Phép chia hết:}\\
$648 : 3 = \mathbf{216}$.\\
-- $6$ chia $3$ được $2$, viết $2$; $2 \times 3 = 6$; $6 - 6 = 0$.\\
-- Hạ $4$; $4$ chia $3$ được $1$, viết $1$; $1 \times 3 = 3$; $4 - 3 = 1$.\\
-- Hạ $8$ được $18$; $18 : 3 = 6$, viết $6$; $6 \times 3 = 18$; $18 - 18 = 0$.\\
\textbf{2. Có chữ số 0 ở thương:}\\
$816 : 4 = \mathbf{204}$.\\
$8 : 4 = 2$; hạ $1$, $1 < 4 \implies$ \textbf{viết 0 vào thương}; hạ $6$ được $16$; $16 : 4 = 4$.
}
\end{tcolorbox}
\end{minipage}
\hfill
\begin{minipage}[t]{0.49\textwidth}\vspace{0pt}
\begin{tcolorbox}[
  enhanced,
  colback=red!4!white,
  colframe=red!70!black,
  arc=2mm, boxrule=0.8pt,
  left=2mm, right=2mm, top=1.8mm, bottom=1.8mm,
  title={\footnotesize\bfseries\color{white}\faExclamationTriangle\ Mẹo tránh bẫy trắc nghiệm}
]
{\footnotesize
\begin{itemize}\setlength{\itemsep}{1.5pt}\setlength{\leftskip}{-2.5mm}
  \item \textbf{Bẫy quên viết chữ số 0 ở thương:}\\
  Ở ví dụ $816 : 4$, rất nhiều bạn hạ $1$ thấy không chia được liền hạ tiếp $6$ thành $16$ rồi ghi thương là $24$! Kết quả đúng phải là $\mathbf{204}$!
  \item \textbf{Mẹo thử lại nhẩm:}\\
  Nếu tính ra thương $24$, nhẩm $24 \times 4 < 100 \ne 816$, phát hiện ngay sai sót!
\end{itemize}
}
\end{tcolorbox}
\end{minipage}


\end{document}
